\documentclass[11pt,a4paper]{article}
\usepackage[margin=2.5cm]{geometry}
\usepackage{amsmath, amssymb}
\usepackage{booktabs}
\usepackage{hyperref}
\usepackage{parskip}
\usepackage{listings}
\usepackage{xcolor}
\usepackage{graphicx}
\usepackage{float}

\lstset{
    basicstyle=\ttfamily\small,
    breaklines=true,
    frame=single,
    backgroundcolor=\color{gray!8},
    keywordstyle=\color{blue},
    commentstyle=\color{green!60!black},
}

\title{\textbf{Reinforcement Learning for Closed-Loop Insulin Regulation}\\
\large Design Evolution and Iterative Fixes}
\author{}
\date{}

\begin{document}
\maketitle

% ─────────────────────────────────────────────────────────────────────────────
\section{Problem Definitions}
% ─────────────────────────────────────────────────────────────────────────────

\subsection{Task}
The goal is to train a reinforcement learning (RL) agent to act as a closed-loop insulin pump
controller for a Type 1 Diabetes (T1D) patient. The agent observes continuous glucose monitor
(CGM) readings and decides how much insulin to deliver every 5 minutes over a 24-hour period
(288 time steps), mimicking an Artificial Pancreas System (APS).

The simulation environment is \textbf{simglucose} (v0.2.11), a Python implementation of the
FDA-accepted UVA/Padova Type 1 Diabetes simulator. All experiments target \texttt{child\#006}
(patient index 15 in the full 30-patient list).

\subsection{State Space}

The state at each time step is a matrix of shape $[\text{seq\_len} \times n_\text{features}]$,
representing a sliding window of recent observations. The feature set is determined by the
\texttt{StateSpace} class from the G2P2C repository.

All three announcement flags default to \texttt{True} in the G2P2C \texttt{Options} parser,
so the active configuration uses $n_\text{features} = 7$:

\begin{enumerate}
    \item \textbf{CGM}: scaled blood glucose $\in [-1,1]$
    \item \textbf{Insulin}: scaled dose delivered $\in [-1,1]$
    \item \textbf{Meal announcement} (\texttt{use\_meal\_announcement}): time remaining
          until the next meal, scaled over $[0, t_\text{meal}]$. Gives the agent advance
          warning before a carbohydrate load arrives.
    \item \textbf{Time of day} (\texttt{use\_tod\_announcement}): step index scaled over
          $[0, 312]$, allowing the agent to learn time-dependent patterns
          (e.g.\ dawn phenomenon).
    \item \textbf{IOB-20}: sum of insulin delivered in the last 20 minutes (4 steps).
    \item \textbf{IOB-60}: sum of insulin delivered in the last 60 minutes (12 steps).
    \item \textbf{IOB-120}: sum of insulin delivered in the last 120 minutes (24 steps).
\end{enumerate}

IOB (Insulin On Board) features 5--7 tell the agent how much active insulin is already
circulating, preventing it from stacking doses and causing hypoglycaemia.

The full state matrix fed to the LSTM has shape $[48 \times 7]$,
where $\text{seq\_len} = \text{feature\_history} = 48$ corresponds to 4 hours of history,
giving $48 \times 7 = 336$ input values per time step.

\subsection{Action Space}

The raw policy output $a \in [-1, 1]$ is converted to a physical insulin dose via the
\textbf{Pump} class using an exponential mapping:

\begin{equation}
    u_\text{pump} = \alpha \cdot \exp\!\bigl((a - 1) \cdot 4\bigr)
    \label{eq:pump}
\end{equation}

where $\alpha = \texttt{action\_scale}$ is a patient-specific scaling factor. This mapping
ensures:
\begin{itemize}
    \item $a = -1 \Rightarrow u_\text{pump} \approx 0$ (near-zero insulin)
    \item $a = 0 \Rightarrow u_\text{pump} = \alpha \cdot e^{-4} \approx 0.018\alpha$
    \item $a = +1 \Rightarrow u_\text{pump} = \alpha$ (maximum dose)
\end{itemize}

For \texttt{child\#006} with basal rate $u_\text{basal} = 0.006735\ \text{U/min}$,
setting $\alpha = 0.37$ centers the action space around the basal dose.

\subsection{Units}

All quantities must be in consistent units throughout the pipeline. Table~\ref{tab:units}
summarises the units used in the G2P2C repository and how they propagate.

\begin{table}[H]
\centering
\begin{tabular}{llll}
\toprule
\textbf{Quantity} & \textbf{Unit} & \textbf{Source} & \textbf{Notes} \\
\midrule
Blood glucose (CGM) & mg/dL & simglucose & Sensor: GuardianRT \\
Insulin dose ($u$)  & U/min & simglucose \texttt{env.step()} & NOT U/h \\
Basal rate          & U/min & \texttt{pumpAction.py} formula & $u_\text{basal} = u2ss \cdot BW / 6000$ \\
Time step           & 5 min & GuardianRT sensor period & 288 steps = 24 h \\
Meal size           & g CHO & simglucose scenario & Carbohydrates in grams \\
\bottomrule
\end{tabular}
\caption{Units used throughout the simulation pipeline.}
\label{tab:units}
\end{table}

The basal rate formula $u_\text{basal} = u2ss \cdot BW / 6000$ converts steady-state insulin
infusion rate ($u2ss$, in pmol/kg/min) and body weight ($BW$, in kg) to U/min. For
\texttt{child\#006}: $u_\text{basal} \approx 0.00673\ \text{U/min}$.

\subsection{Reward Functions}

Two reward functions were used across the project's evolution.

\subsubsection*{Magni Risk Index (repo baseline)}
\begin{equation}
    f(BG) = 1.509\,\bigl(\ln(BG)^{1.084} - 5.381\bigr), \qquad
    r = -\bigl(\text{LBGI} + \text{HBGI}\bigr)
\end{equation}
where $\text{LBGI} = 10 f^2$ when $f < 0$ and $\text{HBGI} = 10 f^2$ when $f > 0$.
The composite reward then normalises this to approximately $[0, 1]$ and applies a hard
penalty of $-15$ for $BG \le 40\ \text{mg/dL}$.

\textbf{Limitation}: The risk index is centred at $BG \approx 112.5\ \text{mg/dL}$ and
penalises hypo and hyper roughly symmetrically. It provides a weak gradient above
$180\ \text{mg/dL}$, making it difficult for the agent to learn post-meal bolusing.

\subsubsection*{Clinical Zone Reward (final version)}
A piecewise-linear reward function aligned with clinical T1D guidelines:

\begin{equation}
r(BG) = \begin{cases}
-15 & BG < 54 \\
-2 - 6\,\dfrac{70 - BG}{16} & 54 \le BG < 70 \\
-0.5\,\dfrac{80 - BG}{10} & 70 \le BG < 80 \\
+1 & 80 \le BG \le 130 \quad \text{(pre-meal target)} \\
1 - 0.5\,\dfrac{BG - 130}{50} & 130 < BG \le 180 \\
0.5 - 1.0\,\dfrac{BG - 180}{40} & 180 < BG \le 220 \quad \text{(post-meal ceiling)} \\
-0.5 - 1.5\,\dfrac{BG - 220}{80} & 220 < BG \le 300 \\
-2 - 3\,\min\!\Bigl(\dfrac{BG-300}{100},1\Bigr) & BG > 300
\end{cases}
\label{eq:clinical_reward}
\end{equation}

This gives the agent a clear gradient to bring glucose into the 80--130 mg/dL range,
and a steeper penalty above 220 mg/dL to encourage post-meal insulin correction.

% ─────────────────────────────────────────────────────────────────────────────
\section{Phase 0 — Original Code}
% ─────────────────────────────────────────────────────────────────────────────

\subsection{Architecture}
The initial agent used a \textbf{Transformer encoder} in \texttt{others.py} and a
\textbf{PPO} training loop in \texttt{agent.py} and \texttt{main.py}. The actor and critic
\emph{shared} the same encoder. Key hyperparameters:

\begin{itemize}
    \item Policy: Transformer $\to$ 3-layer MLP $\to$ Gaussian $(\mu, \sigma)$
    \item PPO: 5 epochs, $\epsilon_\text{clip} = 0.1$
    \item No GAE (plain Monte Carlo returns)
    \item No reward normalisation
    \item Action scale defaulted to 1 (incorrect for child\#006)
\end{itemize}

\subsection{Problems Identified}
\begin{enumerate}
    \item \textbf{Import errors}: \texttt{ModuleNotFoundError} for \texttt{agents}, \texttt{utils},
          etc.\ because the G2P2C repo path was not in \texttt{sys.path}.
    \item \textbf{Patient indexing}: \texttt{patients[15]} failed because the full 30-patient
          list was not loaded. Fixed by using \texttt{get\_patient\_env()} from the repo.
    \item \textbf{Wrong action scale}: Default $\alpha = 1$ places the exponential mapping's
          centre far above child\#006's basal rate, causing the policy to emit excessive insulin
          from the first episode.
    \item \textbf{No basal floor}: The policy could output near-zero insulin (action $\to -1$),
          causing the patient to crash into severe hypoglycaemia.
    \item \textbf{Shared encoder gradient conflict}: The actor and critic shared a single
          encoder; backpropagating the critic loss modified the encoder before the actor
          backward pass, producing a \texttt{RuntimeError: one of the variables needed for
          gradient computation has been modified by an inplace operation}.
    \item \textbf{Missing GAE}: Without Generalised Advantage Estimation, the policy received
          high-variance returns, slowing convergence.
\end{enumerate}

% ─────────────────────────────────────────────────────────────────────────────
\section{Phase 1 — A2C Adaptation from G2P2C Repo}
% ─────────────────────────────────────────────────────────────────────────────

\subsection{Changes Made}
Inspired by the proven G2P2C A2C agent, the architecture was rebuilt to match the repo design:

\begin{itemize}
    \item \textbf{Transformer $\to$ LSTM}: Replaced the Transformer encoder with a single-layer
          LSTM (hidden size 16) in both actor and critic. The LSTM processes the
          $[12 \times n_\text{features}]$ history matrix and outputs a 16-dimensional vector
          from the final hidden state.
    \item \textbf{Independent encoders}: Actor and critic each received their own separate
          \texttt{LSTMEncoder} instance, eliminating the gradient conflict.
    \item \textbf{NormedLinear output heads}: Weight-normalised linear layers (scaled by 0.1)
          replaced standard \texttt{nn.Linear} for $\mu$, $\sigma$, and value output,
          stabilising early training.
    \item \textbf{GAE}: Generalised Advantage Estimation ($\gamma = 0.99$, $\lambda = 0.95$)
          was added, matching the repo's \texttt{compute\_gae()} implementation.
    \item \textbf{Reward normalisation}: Running mean/variance of discounted returns
          (Welford's algorithm) was used to normalise rewards before computing advantages.
    \item \textbf{A2C update}: Single-pass policy gradient
          $\mathcal{L}_\pi = -\mathbb{E}[\log\pi(a|s) \cdot \hat{A}] - \beta H(\pi)$
          with $\beta = 0.001$.
\end{itemize}

\subsection{Remaining Problems}
Although the agent completed full 24-hour episodes without crashing, two issues persisted:
\begin{enumerate}
    \item \textbf{Uniform insulin}: The insulin trace was nearly flat, indicating the policy
          had not learned to bolus in response to meals.
    \item \textbf{Post-meal CGM spikes}: Glucose rose to 360--430 mg/dL after each meal
          because the agent's flat dosing strategy provided no correction bolus.
\end{enumerate}

The root cause was identified as the weak hyperglycaemia signal from \texttt{composite\_reward}:
at $BG = 300\ \text{mg/dL}$, the composite reward is still only $\approx 0.7$ (on a 0--1 scale),
giving the agent little incentive to increase insulin after meals.

% ─────────────────────────────────────────────────────────────────────────────
\section{Phase 2 — Discrete Action Space}
% ─────────────────────────────────────────────────────────────────────────────

\subsection{Motivation}
The continuous exponential action mapping (Eq.~\ref{eq:pump}) creates a sensitive control
surface: small changes in raw output produce large changes in physical insulin dose. This
makes stable learning difficult with a single training worker. A discrete action space
was trialled to simplify exploration.

\subsection{Changes Made}
\begin{itemize}
    \item The Gaussian policy was replaced with a \textbf{Categorical} policy over
          $K = 6$ discrete dose levels defined as multiples of the patient's basal rate:
          \begin{equation}
              \mathcal{A} = \{0.5,\; 1.0,\; 2.0,\; 4.0,\; 8.0,\; 12.0\} \times u_\text{basal}
          \end{equation}
          The policy outputs 6 logits; a \texttt{Categorical} distribution is formed and
          one index is sampled per step.
    \item The Pump exponential mapping was removed entirely. Each action index maps
          directly to a fixed physical dose, so $u_\text{pump} = \texttt{DOSE\_MULTIPLIERS}[i]
          \times u_\text{basal}$.
    \item Log probability computation changed from \texttt{Normal.log\_prob} to
          \texttt{Categorical.log\_prob}; the PPO ratio and clipping remain identical.
    \item The training loop was simplified to \textbf{one PPO update per episode}
          (288 steps), giving 2,800 episodes total ($\approx$806,400 transitions).
    \item \texttt{select\_action} now returns \texttt{(pump\_act, log\_prob, action\_index)};
          the agent's \texttt{update()} takes integer indices as a \texttt{LongTensor}
          instead of floating-point raw actions.
\end{itemize}

\subsection{Result and Problems}
The insulin trace showed slightly more variation than Phase 1, but CGM still spiked
post-meal and the policy defaulted to the lowest non-zero dose level most of the time.
The discrete agent did not meaningfully outperform the continuous one at this training
budget. The approach was abandoned in favour of returning to continuous PPO with
a cleaner architecture.

% ─────────────────────────────────────────────────────────────────────────────
\section{Phase 3 — Plain Vanilla PPO (Continuous)}
% ─────────────────────────────────────────────────────────────────────────────

\subsection{Changes Made}
The agent was fully reverted to \textbf{continuous PPO} matching the G2P2C PPO agent
(\texttt{agents/ppo/}):

\begin{itemize}
    \item \textbf{PPO clipped surrogate loss}:
    \begin{equation}
        \mathcal{L}_\pi = -\mathbb{E}\Bigl[\min\bigl(r_t \hat{A}_t,\;
        \text{clip}(r_t, 1-\epsilon, 1+\epsilon)\,\hat{A}_t\bigr)\Bigr]
        - \beta H(\pi), \quad \epsilon = 0.1
    \end{equation}
    where $r_t = \pi_\theta(a|s) / \pi_{\theta_\text{old}}(a|s)$ is the importance
    sampling ratio.
    \item \textbf{5 update epochs} per rollout batch (repo: \texttt{n\_pi\_epochs = 5}).
    \item \textbf{Separate optimisers} for actor and critic with $\text{lr} = 3 \times 10^{-4}$.
    \item \textbf{Basal floor}: After the Pump mapping, $u_\text{pump}$ is clipped to
          $\max(u_\text{pump},\; u_\text{basal})$ to prevent the policy from collapsing to
          near-zero insulin output.
    \item \textbf{Gradient clipping}: Gradient norm clipped to 20 for both networks.
    \item \textbf{Training budget}: 2,800 episodes $\times$ 288 steps $\approx$ 806,400
          transitions, matching the repo's $\sim$800k interaction budget. One PPO update
          is performed per episode.
\end{itemize}

\subsection{Result}
The agent successfully learned basal regulation: CGM remained in approximately
110--180 mg/dL during fasting periods. Post-meal behaviour improved — CGM peaked around
320--350 mg/dL and returned toward target after $\sim$40 steps. However, the glucose
rarely entered the pre-meal target of 80--130 mg/dL, and one meal event caused the
insulin to drop (rather than rise), suggesting the policy had not fully learned the
causal link between meals and required bolus dosing.

% ─────────────────────────────────────────────────────────────────────────────
\section{Phase 4 — Clinical Zone Reward}
% ─────────────────────────────────────────────────────────────────────────────

\subsection{Motivation}
The Magni risk index is symmetric around 112.5 mg/dL and provides a weak gradient
above 180 mg/dL. Clinical T1D guidelines define clear target zones with distinct
urgency levels that are not captured by the risk index. The \texttt{composite\_reward}
was replaced entirely with the zone-based \texttt{clinical\_reward} (Eq.~\ref{eq:clinical_reward}).

\subsection{Changes Made}
\begin{itemize}
    \item \textbf{Removed \texttt{composite\_reward}} import and call entirely.
    \item \textbf{Added \texttt{clinical\_reward}} defined inline in \texttt{main.py},
          returning a scalar for a single CGM reading.
    \item The piecewise structure provides:
    \begin{itemize}
        \item A \textbf{maximum reward of 1.0} for $BG \in [80, 130]\ \text{mg/dL}$
              (pre-meal clinical target).
        \item A \textbf{soft gradient} from 1.0 to 0.5 for $BG \in [130, 180]$
              (acceptable range).
        \item A \textbf{moderate penalty} transitioning to negative values above 180 mg/dL,
              with steeper penalties above 220 mg/dL (post-meal realistic ceiling).
        \item A \textbf{hard cliff of $-15$} below 54 mg/dL, matching the repo's
              treatment of severe hypoglycaemia.
    \end{itemize}
\end{itemize}

\subsection{Expected Improvement}
By providing a clear, continuously differentiable (except at zone boundaries) reward
signal that peaks in the clinical target range and strongly penalises post-meal
hyperglycaemia, the agent should receive a much stronger learning signal to issue
correction boluses after meals and maintain tighter glucose control during fasting.

% ─────────────────────────────────────────────────────────────────────────────
\section{Summary of Evolution}
% ─────────────────────────────────────────────────────────────────────────────

\begin{table}[H]
\centering
\renewcommand{\arraystretch}{1.3}
\begin{tabular}{lllll}
\toprule
\textbf{Phase} & \textbf{Policy} & \textbf{Encoder} & \textbf{Update} & \textbf{Reward} \\
\midrule
0 — Original     & Gaussian (PPO)    & Transformer (shared) & PPO 5ep, no GAE  & composite (risk) \\
1 — A2C Adapt    & Gaussian (A2C)    & LSTM (separate)      & A2C 1ep, GAE     & composite (risk) \\
2 — Discrete     & Categorical       & LSTM (separate)      & PPO 5ep, GAE     & composite (risk) \\
3 — Vanilla PPO  & Gaussian (PPO)    & LSTM (separate)      & PPO 5ep, GAE     & composite (risk) \\
4 — Clinical     & Gaussian (PPO)    & LSTM (separate)      & PPO 5ep, GAE     & clinical zones   \\
\bottomrule
\end{tabular}
\caption{Architecture and reward function at each phase.}
\label{tab:summary}
\end{table}

\end{document}
